\section{Dữ liệu và thiết kế đánh giá}

Chuỗi gốc gồm \DatasetN{} quan sát ngày, từ \DateMin{} đến \DateMax{}, không thiếu
ngày, không trùng ngày và không có giá trị lưu lượng thiếu. Lưu lượng trung bình là
\QMean{} m³/s, độ lệch chuẩn \QStd{} m³/s, nhỏ nhất \QMin{} m³/s và lớn nhất
\QMax{} m³/s. Q95 tính trên toàn bộ chuỗi là \QHigh{} m³/s; có \NDaysHigh{} ngày vượt
ngưỡng (\ShareHigh\% tương đương 5\%), tạo thành \NEvents{} đợt liên tiếp. Z tương
đương ngưỡng này là \ZThreshold, không phải 1,645; vì vậy không dùng z $>$ 3 để thay
cho định nghĩa Top 5\%.

Dữ liệu được chia theo thời gian, không xáo trộn. Train dùng 2010--2019, Validation
dùng 2020--2021 và Test dùng 2022--2023. Mỗi ngày gốc tạo 7 cặp dự báo D+1 đến D+7;
cửa sổ mục tiêu không vượt khỏi tập tương ứng. Bảng~\ref{tab:split} tóm tắt số ngày
gốc. Sau khi chọn mô hình, toàn bộ 2010--2021 được dùng để refit.

\begin{table}[H]
\centering
\caption{Chia tập theo thời gian}
\label{tab:split}
\small
\begin{tabular}{llr}
\toprule
Tập & Khoảng ngày mục tiêu & Số gốc dự báo \\
\midrule
Train & 01/01/2010–31/12/2019 & 3.615 \\
Validation & 01/01/2020–31/12/2021 & 725 \\
Refit & 01/01/2010–31/12/2021 & 4.346 \\
Test & 01/01/2022–31/12/2023 & 724 \\
\bottomrule
\end{tabular}

\end{table}

Ba mô hình được so sánh. Naive persistence dùng $\hat y(t+h)=y(t)$. Seasonal Naive
365 dùng giá trị cùng ngày của năm trước. LightGBM direct 7 dùng lag 1, 2, 3, 7, 14,
28 ngày; trung bình và độ lệch chuẩn trượt 7, 14, 30, 90 ngày; biến lịch sin/cos của
ngày trong năm và tháng; cùng biến lịch của ngày mục tiêu. Mục tiêu của LightGBM là
chênh lệch $q(t+h)-q(t)$, sau đó cộng lại vào lưu lượng ngày gốc; đây vẫn là mô hình
direct cho từng horizon. Siêu tham số được chọn trên Validation, Test chỉ dùng một lần
để báo cáo.

Với mỗi horizon, RMSE, MAE và NSE được tính từng ngày gốc; mean RMSE/MAE/NSE lấy
trung bình trên 7 horizon, còn pooled tính trên toàn bộ cặp dự báo. NSE = 1 là hoàn
hảo, NSE = 0 tương đương dự báo bằng trung bình và NSE âm là kém hơn mức trung bình.
